\textit{This  chapter discusses the societal and economic consequences, or lack thereof the project contains, while discussing sustainability and environmental sides.}

\section{Societal and Ethical Impact}


In communication with the customer on societal and ethical impact, key insights were given. Better visualization of AIS anomalies, is a tool that would transform how current work is done, rather than remove employment opportunities, leading to specialized workers, with access to more effective tools. As such the workplace impact of a fully developed software would not have a negative impacts on workers, neither would the use of better tools lead to loss of employment. As the system would be used to analyze data from large vessels, small vessels would not be affected by this. As a result no surveilance of individual persons would be committed if the tool was connected to real data. A fact of this is that boats are required to use AIS transponders by law, leading to the system not being a new form of surveillance on sea vessels. Another point raised in talks with customer is around the scenario if this could be used to stop illegal smuggling or immigration, but the vessel type being large, means the likelihood of this is low, as large vessels are less suited for illegal tasks like smuggling. 


\section{Economic}


The prototype itself, proves that feature testing can be used to see value, or lack thereof in new features and systems. This can reduce financial risk through not implementing wasted features, leading to less negative consequences economically. Also the use of open source tools cut down on licensing and other costs, leading to the prototype being a financially viable way of testing new systems and features. 

\section{Sustainability}

The sustainability impact of a software development project is generally low. The project itself can reduce wasted effort implementing features not needed. THE system itself can be extended and is maintainable for other developers to work on, leading to a scalable and maintainable system. 


\section{Environment}

The software not requiring physical components aside from computers, leads to very low emissions. And only electricity being required to run the software. Improved AIS detection could lead to action being taken sooner on faulty boats, leading to reduced environmental impact. The software itself can then not have a large amount of environmental impact. 